\section{Задание 1}
В языке Prolog аргументами вопроса могут быть атомы --- конкретные значения 
с маленькой буквы или переменные --- неизвестные значения с большой буквы.

\subsection{Вопрос с двумя атомами}
Вопрос~\ref{lst:t1-q1} из примера содержит два атома и проверяет правда ли существует
связь \code{parent} между \code{don} и \code{mike}. Ответом на данный вопрос является 
истина или ложь.
\begin{listing}[Вопрос с двумя атомами]{lst:t1-q1}{Prolog}
?- parent(don, mike).

true.
\end{listing}

\subsection{Вопрос с атомом и переменной}
Если один из аргументов вопроса является переменной, то система проверяет, существует
ли такое значение этой переменной, для которого удовлетворяется отношение.
Например, ответом на вопрос~\ref{lst:t1-q2} являются все значения \code{Child}, для которых удовлетворяется 
отношение \code{parent}, то есть все дети \code{don}.
\begin{listing}[Вопрос с атомом и переменной]{lst:t1-q2}{Prolog}
?- parent(don, Child).

Child = randy 
Child = mike 
Child = anne
\end{listing}

Ответом на вопрос~\ref{lst:t1-q3} являются все значения \code{Parent}, для которых удовлетворяется 
отношение \code{parent}, то есть все родители \code{randy}.
\begin{listing}[Вопрос с атомом и переменной]{lst:t1-q3}{Prolog}
?- parent(Parent, randy).

Parent = don
Parent = rosie
\end{listing}

\subsection{Вопрос с двумя переменными}
Если оба аргумента вопроса являются переменными, то система ищет все такие значения 
переменных, для которых удовлетворяется отношение. Таким образом ответом на 
вопрос~\ref{lst:t1-q4} являются все значения \code{Parent} и \code{Child}, для которых
выполняется отношение \code{parent}.

\begin{listing}[Вопрос с двумя переменными]{lst:t1-q4}{Prolog}
?- parent(Parent, Child).

Child = randy,
Parent = don
Child = mike,
Parent = don
Child = anne,
Parent = don
Child = randy,
Parent = rosie
Child = mike,
Parent = rosie
Child = anne,
Parent = rosie
\end{listing}

\subsection{Вопрос с несуществующим атомом}
Если хотя бы один из аргументов вопроса является атомом, который не присутствует в базе знаний,
то ответом на такой вопрос является ложь. Пример представлен в листинге~\ref{lst:t1-q5}.
\begin{listing}[Вопрос с несуществующим атомом]{lst:t1-q5}{Prolog}
?- parent(don, john).

false.
\end{listing}

\section{Задание 2}
\subsection{Исходный код}
Исходный код программы <<Телефонный справочник и автомобили>> приведен в листинге~\ref{lst:t2-prog}.
\begin{listing}[Исходный код программы <<Телефонный справочник и автомобили>>]{lst:t2-prog}{Prolog}
contact(ivanov, '111-11-11', address(moscow, lenina, 10, 5)).
contact(ivanov, '222-22-22', address(moscow, lenina, 10, 5)).

contact(ivanov, '333-33-33', address(kazan, pushkina, 15, 42)).

contact(petrov, '444-44-44', address(moscow, tverskaya, 1, 12)).
contact(petrov, '444-44-45', address(moscow, tverskaya, 1, 12)).

contact(sidorov, '555-55-55', address(spb, nevsky, 100, 1)).


car(ivanov, bmw, black, 50000, 'A111AA177').
car(ivanov, bmw, black, 100000, 'A001MP97').
car(ivanov, lada, white, 5000, 'B222BB97').

car(petrov, bmw, black, 55000, 'X999XX77').
car(petrov, porsche, gray, 55000, 'E999KX97').

car(sidorov, audi, red, 40000, 'O000OO78').
car(sidorov, audi, blue, 40000, 'O001OO78').
car(sidorov, bmw, black, 400000, 'O777OO178').


search(LastName, Phone, City, Street, House, Flat, Brand, Color, Price, RegNum) :-
    car(LastName, Brand, Color, Price, RegNum),
    contact(LastName, Phone, address(City, Street, House, Flat)).
\end{listing}

\subsection{Что собой представляет программа на Prolog?}
Программа представляет собой базу знаний и вопрос. Программа на Prolog описывает факты 
предметной области и логические правила связей между ними.

\subsection{Какова структура программы на Prolog?}
Программа на языке SWI-Prolog не имеет жестко заданных разделов с объявлением переменных и не 
содержит точки входа. Структура разделена на два смысловых блока:
\begin{itemize}
    \item факты:
    \begin{itemize}
        \item отношение \code{contact} описывает абонента, для группировки адресных данных 
        используется составной терм \code{address(Город, Улица, Дом, Квартира)};
        \item отношение \code{car} описывает транспортное средство владельца;
    \end{itemize}
    \item правила:
    \begin{itemize}
        \item конъюнктивное правило \code{search} связывает два разных отношения.
    \end{itemize}
\end{itemize}

\subsection{Как реализуется программа?}
Программа реализуется с использованием единственной основной конструкции языка --- термов.
\subsubsection{Факты}
Справочник абонентов и автомобилей реализован в виде основных фактов --- безусловных утверждений о том, 
что между объектами существует связь.

\subsubsection{Составные термы}
Для учета сложной структуры данных применяется составной терм \code{address(...)}, 
где \code{address} выступает в роли функтора, объединяющего аргументы в единый объект.

\subsubsection{Правила}
Логика поиска реализована через правило \code{search}, содержащее:
\begin{itemize}
    \item заголовок --- составной терм, фиксирующий отношение между исходными и искомыми объектами;
    \item тело --- набор условий истинности знания.
\end{itemize}

\subsubsection{Использование переменных}
Для передачи значений в пространстве и во времени используется именованная переменная. Она позволяет 
в обобщенной форме зафиксировать связь для любого элемента множества.
Для сокрытия лишней информации используются анонимные переменные, который не связываются ни с каким 
значением.

\subsection{Как формируются результаты работы программы}
Результаты работы программы формируются путем получения ответов системы на заданный 
целевой вопрос. Процесс доказательства истинности запроса базируется на следующих механизмах.

\subsubsection{Применение правил вывода и унификация}
Так как факты в базе знаний являются основными, а целевой вопрос и правило 
\code{search} содержат переменные, система использует правило обобщения факта. 
Сопоставление объектов в процессе вывода реализуется через операцию унификации. 
Два терма считаются унифицируемыми, если выполняется одно из условий:
\begin{itemize}
    \item $T_1$ и $T_2$ --- идентичные константы;
    \item $T_1$ --- неконкретизированная переменная, а $T_2$ --- константа 
    или составной терм, не содержащий $T_1$ в качестве аргумента. В этом случае 
    $T_1$ конкретизируется значением $T_2$;
    \item $T_1$ и $T_2$ --- неконкретизированные переменные, которые в результате 
    сопоставления становятся сцепленными;
    \item $T_1$ и $T_2$ --- составные термы с одинаковыми функторами и 
    арностью, при условии успешной унификации каждой пары их соответствующих аргументов.
\end{itemize}

\subsubsection{Формирование подстановок}
В ходе работы алгоритма унификации система находит подстановку 
$\theta$ --- множество пар вида $\{x_i = t_i\}$, где каждая неизвестная 
переменная $x_i$ заменяется на соответствующий конкретный терм $t_i$ из 
базы знаний. Применение подстановки $\theta$ к исходному терму позволяет 
получить его пример частный случай, который является истинным в рамках данной системы знаний.

\subsubsection{Поиск подтверждения и механизм реконкретизации}
Процесс формирования ответа носит итерационный характер. При вызове 
конъюнктивного правила система последовательно доказывает истинность 
каждого терма в его теле. В процессе поиска именованные переменные конкретизируются. 

Если на определенном этапе доказательства очередное условие не может быть 
выполнено, система использует механизм отката backtracking. При этом происходит 
реконкретизация переменных, и система возвращается к предыдущему шагу для поиска 
альтернативных подстановок. Работа программы продолжается до тех пор, пока не 
будет найдена полная подстановка, удовлетворяющая всем условиям правила, либо 
пока не будет исчерпано всё пространство поиска в базе знаний.
